\section{楕円ブロックの明示模型と有限係数の主定理}\label{sec:explicit-block-results}

\[
 S_h=\CC\langle x,y,z\rangle/
 ([x,y]-hz^2,\ [y,z]-hx^2,\ [z,x]-hy^2)
\]
と置き、\(T=\sym(xyz)\) を六つの順列の平均、
\(G_h=x^3+y^3+z^3-2h^2T\) とする。
中心変数 \(u_1,\ldots,u_m\) と斉次三次式 \(P(u)\) に対して
\[
 B_{h,P}=S_h[u_1,\ldots,u_m]/(G_h+P(u))
\]
を考える。\(P=\sum_i u_i^3\) の場合を \(B_{h,m}\) と略記する。

\begin{theorem}[一ブロックと任意の中心三次式の可換モデル]\label{thm:global}
 \(m\ge2\)、\(n=m+1\)、\(h\in\CC\)、\(h^6\ne-27\) とする。
 \(P\in\CC[u_1,\ldots,u_m]_3\) のアフィン臨界点が原点だけであると仮定する。
 \(B_{h,P}\) の残余dg圏とrotationは、滑らかなcubic
 \[
 Y_{h,P}=V\!\left(x^3+y^3+z^3+h^2xyz+P(u)\right)
 \subset\PP^{m+2}
 \]
 の \((\Ku(Y_{h,P}),\rho_{Y_{h,P}})\) と同値である。
 この \(Y_{h,P}\) は、同じrotation付き残余圏を持つ滑らかな可換cubicとして
 射影同型を除いて一意である。
\end{theorem}

この定理はsixfoldでは \(m=5\)、sevenfoldでは \(m=6\) を含む。
中心変数どうしの任意の三次混合項を許すが、一つのSklyaninブロックと中心部分が
分離されているという仮定を保持する。\(P=\sum u_i^3\) はその特殊例である。
複数の非可換ブロックについても、対応する可換cubicの残余圏からの完全忠実埋込みを得る
（命題\ref{prop:multiblock-embedding}）。次の部分的な有限位数条件の下では、その像が全体を生成する。

\begin{theorem}[一つの任意位数ブロックを許す多ブロック可換モデル]\label{thm:multiblock-model}
 \(r\ge2\)、\(m\ge0\)、\(n=3r+m-2\) とする。
 各 \(j\) について \(h_j^6\ne-27\) とし、楕円曲線
 \[
 E_{h_j}=V(x_j^3+y_j^3+z_j^3+h_j^2x_jy_jz_j),\qquad
 o_j=[1:-1:0],\quad p_j=[1:-1:-h_j]
 \]
 上の点 \(p_j\) は、高々一つを除いて有限位数であると仮定する。
 すなわち因子を並べ替えた後、\(j\ge2\) にのみ有限位数を仮定し、\(p_1\) の位数には条件を課さない。
 \(m>0\) の場合、\(P\in\CC[u_1,\ldots,u_m]_3\) は原点だけに臨界点を持つとし、
 \(m=0\) の場合は中心部分を省く。
 異なるSklyaninブロックの元を可換とした代数
 \[
 B_{\mathbf h,P}=
 (S_{h_1}\otimes\cdots\otimes S_{h_r})[u_1,\ldots,u_m]
 \big/\left(\sum_jG_{h_j}+P\right)
 \]
 のrotation付き残余dg圏は、滑らかなcubic
 \[
 Y_{\mathbf h,P}=V\!\left(
 \sum_{j=1}^r(x_j^3+y_j^3+z_j^3+h_j^2x_jy_jz_j)+P(u)
 \right)\subset\PP^{3r+m-1}
 \]
 の \((\Ku(Y_{\mathbf h,P}),\rho_{Y_{\mathbf h,P}})\) と同値である。
 この可換モデルは射影同型を除いて一意である。
\end{theorem}

ここで「有限パラメータ」は形式変数でなく実際の複素数であることを意味し、
楕円曲線上の点が「有限位数」であることとは異なる。
定理は一つの無限位数ブロックを許す。
具体的に \(h_1=2\) は無限位数（補題\ref{lem:order-infinite}）、
\(h_j^6=5\ (j\ge2)\) は位数五（補題\ref{lem:order-five}）となる。
従って二ブロックのsixfold、三ブロックのsevenfoldにも、真に無限位数の因子を含む例が得られる。
二つ以上の無限位数ブロックを任意に許す主張は、本定理には含まれない。
証明は第\ref{sec:finite-center}節に置く。

\begin{proposition}[例外パラメータのhomological singularity]\label{prop:bad-intro}
 \(m\ge0\) とする。Fermat中心部分の族 \(B_{h,m}\) で \(h^6=-27\) とすると、\(S_h[u]\) はAS-regularで
 \(G_h+\sum u_i^3\) は中心正則元であるが、そのgraded matrix factorization圏は
 properではなく、さらに
 \[
 \operatorname{hdim}(\qgr B_{h,m})=\infty
 \]
 である。ここでhomological dimensionは、非零のExtが存在する次数の上限である。
 従って、滑らかな射影cubicのKuznetsov成分とは同値にならない。
\end{proposition}

これはAS-regularな周囲、三次の中心正則元、期待されるHilbert series、
および \(\rho^3\simeq[2]\) だけでは、大域的な可換モデルの存在を保証しないことを示す。
本稿の可換点からのPoisson量子化族に限定した、良い有限パラメータ全体での存在問題は別に残る。
一般のAS-regular中心三次商については、smooth properかつCalabi--Yauな例にも
定理\ref{thm:zhang-cy-fourfold-counterexample}の反例がある。

